\subsection{射影线性簇的空间}

给定一个除环 K，集 \(\mathbf{P}_{n,p}(\mathbf{K})\) ——\(p \geqslant 0\) 维射影线性簇（左射影空间 \(\mathbf{P}_{n}(\mathbf{K})\) 的）组成的集——显然与 \(p + 1\) 维向量子空间（左向量空间 \(K_{s}^{n+1}\) 的）组成的集一一对应。令 \(L_{n+1, p+1}(\mathbf{K})\) 表示自由系 \((x_{k})_{1 \leqslant k \leqslant p+1}\) ——\(p + 1\) 个向量（\(K_{s}^{n+1}\) 中的）组成的自由系——所成的集；那么集 \(\mathbf{P}_{n,p}(\mathbf{K})\)

仍然与 \(\mathbf{L}_{n+1,p+1}(\mathbf{K})\) 关于下述等价关系 \(\Delta_{n,p}(\mathbf{K})\) 的商一一对应：“\((\mathbf{x}^{k})\) 与 \((\mathbf{y}^{k})\) 生成同一个 \(p+1\) 维向量子空间（\(K_{s}^{n+1}\) 的）”。在下文中我们将把 \(\mathbf{P}_{n,p}(\mathbf{K})\) 与这个商集等同。另一方面，若把每个自由系 \((\mathbf{x}_{k})\) ——\(p+1\) 个向量（\(K_{s}^{n+1}\) 中的）组成的——对应于矩阵 X ——\(p+1\) 行 \(n+1\) 列的，以 \(x_{k}\) 为第 k 行（ \(1\leqslant k\leqslant p+1\) ）——，我们就得到 \(L_{n+1,p+1}(\mathbf{K})\) 与所有 \(p+1\) 行 \(n+1\) 列、秩为 \(p+1\) 的矩阵组成的集之间的一一对应；我们将把 \(L_{n+1,p+1}(\mathbf{K})\) 与这个矩阵集等同，而两个矩阵 X, Y 之间的关系 \(\Delta_{n,p}(\mathbf{K})\) 就是：“存在 \(p+1\) 阶非奇异方矩阵 T 使得 Y = T.X”。

在下文中我们取 K 为域 R，并在上述记号中省略字母 K。我们可以用一种推广实射影空间拓扑定义的办法来定义 \(P_{n,p}\) 上的拓扑。即，\(L_{n+1, p+1}\) 含于空间 \(M_{p+1, n+1}\) ——以实数为元素的 \(p + 1\) 行 \(n + 1\) 列的一切矩阵组成的空间——；我们给 \(L_{n+1, p+1}\) 赋以这个矩阵空间的拓扑所诱导的拓扑（§ 1，第 6 小节）。

定义 2. 空间 \(\mathbf{P}_{n,p}\) ——拓扑空间 \(\mathbf{L}_{n+1,p+1}\) 关于等价关系 \(\Delta_{n,p}\) 的商——称为 \(p\geqslant0\) 维射影线性簇（在实射影空间 \(\mathbf{P}_{n}\) 中的）的空间。

我们将使用下述记号：给定一个 \(p + 1\) 行 \(n + 1\) 列的矩阵 X，以及任一严格递增序列

\[
\sigma = (i _ {1}, \dots , i _ {p + 1})
\]

的 \(p + 1\) 个指标（区间 \([0, n]\) ，属于 \(\mathbf{N}\) ），我们用 \(X_{\sigma}\) 表示 \(X\) 的由指标为 \(i_1, i_2, \ldots, i_{p+1}\) 的列组成的方子矩阵。我们用 \(\mathbf{A}_{\sigma}\) 表示 \(\mathbf{L}_{n+1, p+1}\) 中由这些矩阵 \(X\) ——使 \(X_{\sigma}\) 非奇异的——组成的子集。由 § 1 第 6 小节的命题 6，\(\mathbf{A}_{\sigma}\) 是 \(\mathbf{M}_{p+1, n+1}\) 中的稠密开集，而且函数 \(X \to X_{\sigma}^{-1}\) 在 \(\mathbf{A}_{\sigma}\) 上连续。

集 \(A_{\sigma}\) 的一个几何解释如下：设 \(E_{\sigma}\) 是 \(R^{n+1}\) 中由典范基的向量 \(e_{i}\) （使 \(i \in \sigma\) 的）生成的向量子空间，设 \(E_{\sigma}'\) 是由 \(e_{i}\) （使 \(i \notin \sigma\) 的）生成的余子空间；那么说矩阵 x 属于 \(A_{\sigma}\) ，意思是其在 \(E_{\sigma}\) 上的投影——它的 \(p + 1\) 行（即 \(x_{k}\)）的投影——构成自由系，或者说，由诸 \(x_{k}\) 生成的向量子空间是 \(E_{\sigma}'\) 的余子空间（或者说它与 \(E_{\sigma}'\) 的交只由 o 组成）。

命题 6. 空间 \(\mathbf{P}_{n,p}\) 是豪斯多夫的。

我们首先证明关系 \(\Delta_{n,p}\) 是开的。若 U 是 \(L_{n+1,p+1}\) 中的开集，要使 U 关于 \(\Delta_{n,p}\) 饱和，须取 U 在映射 \(X \to T \cdot X\) 下的象的并，这里 T 取遍 \(p + 1\) 阶非奇异方矩阵的集；由于这些映射中每一个都是双连续的，所有这些象都是开集，因此它们的并也是开集。

由第 I 章，§ 8，第 3 小节的命题 8，只要证明 \(L_{n+1,p+1} \times L_{n+1,p+1}\) 的由 \(\Delta_{n,p}\) 定义的子集 N 是闭的，证明即告完成。设 \((X,\mathcal{Y})\) 是乘积空间中属于 N 的闭包的一点，设 \(\sigma\) 是使 \(X_{\sigma}\) 非奇异的指标序列：由于 \(A_{\sigma}\) 是开的，存在 \((X,\mathcal{Y})\) 的邻域 V，使得对每对 \((X',\mathcal{Y}') \in \mathrm{N} \cap \mathrm{V}\) ，矩阵 \(X'_{\sigma}\) 都是非奇异的；当 \((X',\mathcal{Y}')\) 保持属于 N 而趋于 \((X,\mathcal{Y})\) 时，矩阵 \(Y'_{\sigma}X_{\sigma}^{\prime -1}\) 就趋于 \(T = Y_{\sigma}X_{\sigma}^{-1}\) ；由于我们有 \(Y' = (Y'_{\sigma}X_{\sigma}^{\prime -1})X'\) ，取极限便见 Y = T.X，证明完成。

命题 7. 空间 \(\mathbf{P}_{n,p}\) 是紧的。

只需证明存在 \(L_{n+1, p+1}\) 的一个紧子空间，它与模 \(\Delta_{n,p}\) 的每个等价类至少交于一点；因为这时 \(P_{n,p}\) 是这个子空间在 \(L_{n+1, p+1}\) 到 \(P_{n,p}\) 上的典范映射下的象，因此是紧的（第 I 章，§ 9，第 4 小节，定理 2）。

令 \(\mathbf{V}_{n+1, p+1}\) 是 \(\mathbf{L}_{n+1, p+1}\) 的这样的子空间：其元素是系 \((\boldsymbol{x}_k)\) ——由 \(p + 1\) 个向量组成的——，这些向量构成它们所生成的向量子空间的欧氏标准正交基，也就是说满足 \((x_h | x_k) = 0\) （每当 \(h \neq k\) ），且 \((x_h | x_h) = 1\) （对 \(1 \leqslant h \leqslant p + 1\) ）。每个 \(p + 1\) 维向量子空间（\(\mathbf{R}^{n+1}\) 的）都有这样一个基，因此模 \(\Delta_{n,p}\) 的每个类都与 \(\mathbf{V}_{n+1, p+1}\) 相交。另一方面，矩阵 \(X = (x_{ij})\) ——\(\mathbf{V}_{n+1, p+1}\) 的——由下述关系定义

\[
\sum_ {j = 0} ^ {n} x _ {i j} ^ {2} = 1 (1 \leqslant i \leqslant p + 1),
\]

\[
\sum_ {j = 0} ^ {n} x _ {i j} x _ {k j} = 0 \quad (i \neq k);
\]

因此它们组成 \(M_{p+1, n+1}\) 中的一个闭集；又由于这些关系蕴含 \(|x_{ij}| \leqslant 1\) （对每对指标 \((i, j)\) ），这个集是有界的，从而是紧的。

命题 8. \(\mathbf{P}_{n,p}\) 是连通的且局部连通的，并且每个点都有一个同胚于 \(\mathbf{R}^{(p+1)(n-p)}\) 的开邻域。

对每个（严格递增的）指标序列 \(\sigma\) ，集 \(\mathbf{A}_{\sigma}\) 在 \(\mathbf{L}_{n+1, p+1}\) 中是开的且关于 \(\Delta_{n,p}\) 饱和；它的典范象

\(C_{\sigma}\) （\(P_{n,p}\) 中的）因此是一个开集，同胚于 \(A_{\sigma}\) 关于等价关系 \(\Theta_{\sigma}\) （在 \(A_{\sigma}\) 上由 \(\Delta_{n,p}\) 诱导的）的商（引理 1）。

令 \(\mathbf{B}_{\sigma}\) 是 \(\mathbf{A}_{\sigma}\) 中由矩阵 \(X\) ——使 \(X_{\sigma}\) 为 \(p + 1\) 阶单位矩阵的——组成的子集；\(X\) 中除 \(X_{\sigma}\) 的元素以外的元素这时是任意的，因此 \(\mathbf{B}_{\sigma}\) 同胚于空间 \(\mathbf{R}^{(p + 1)(n - p)}\) 。对每个矩阵 \(X \in \mathbf{A}_{\sigma}\) ，使它对应于矩阵 \(\Upsilon = X_{\sigma}^{-1}X\) ，它属于 \(\mathbf{B}_{\sigma}\) ；这样我们就定义了连续映射 \(g\) ——\(\mathbf{A}_{\sigma}\) 到 \(\mathbf{B}_{\sigma}\) 上的——，使 \(g(X)\) 是 \(\mathbf{B}_{\sigma}\) 中与 \(X\) 模 \(\Theta_{\sigma}\) 同余的唯一的矩阵。由此推出 \(\mathbf{B}_{\sigma}\) 同胚于 \(\mathbf{A}_{\sigma}/\Theta_{\sigma}\) （引理 2），从而同胚于 \(\mathbf{C}_{\sigma}\) 。

集 \(C_{\sigma}\) 因此是连通的。由于 \(A_{\sigma}\) 在 \(L_{n+1, p+1}\) 中稠密，\(C_{\sigma}\) 在 \(P_{n, p}\) 中稠密，因此 \(P_{n, p}\) 也是连通的（第 I 章，§ 11，第 1 小节，命题 1）。另一方面，\(P_{n, p}\) 的每个点都属于 \(C_{\sigma}\) （对至少一个指标序列 \(\sigma\) ），从而有一个同胚于 \(\mathbf{R}^{(p+1)(n-p)}\) 的开邻域。

矩阵 \(Y=g(X)\) 可以解释如下：为简单起见，设序列 \(\sigma\) 由 \(p+1\) 个指标 n-p, \(n-p+1\) , \(\ldots\) , n 组成，并用 \(a_{ij}\) ( \(1\leqslant i\leqslant p+1\) , \(o\leqslant j\leqslant n-p-1\) ) 表示 y 的前 n-p 列的元素；那么由 x 的行生成的 \(R^{n+1}\) 的向量子空间就是由方程

\[
x _ {j} = \sum_ {i = 1} ^ {p + 1} a _ {i j} x _ {n - p + i - 1} \quad (0 \leqslant j \leqslant n - p - 1).
\]

